\section{Kesimpulan}
Milestone 3 berhasil menambahkan tahap semantic analysis pada Arion Compiler. Program dapat menerima source code, token stream, atau parse tree, kemudian menghasilkan parse tree, AST, symbol table, dan decorated AST. Implementasi AST membuat representasi program menjadi lebih ringkas dibanding parse tree, sedangkan semantic analyzer menambahkan informasi tipe, scope, block, dan referensi symbol table.

Symbol table \texttt{tab}, \texttt{btab}, dan \texttt{atab} telah diimplementasikan untuk menyimpan identifier, block, dan array. Program juga telah mendukung predefined identifier, registrasi deklarasi, lookup identifier, pengecekan redeklarasi, type checking, assignment compatibility, procedure/function call, array access, record field access, serta validasi kondisi control flow.

Berdasarkan pengujian, program berhasil menangani kasus dasar, ekspresi, control flow, structured type, dan subprogram. Program juga mampu melaporkan beberapa semantic error secara informatif tanpa crash.

\section{Saran}
Pengembangan berikutnya dapat menambahkan analisis data-flow yang lebih lengkap untuk mendeteksi penggunaan variabel sebelum inisialisasi pada seluruh jalur eksekusi. Selain itu, pengecekan string length, kompatibilitas record anonymous, dan validasi cakupan nilai subrange pada ekspresi non-literal dapat diperluas agar semakin dekat dengan bahasa Pascal-like yang lengkap.

Untuk tahap selanjutnya, decorated AST yang sudah tersedia dapat dimanfaatkan sebagai dasar intermediate representation atau code generation. Pengujian juga dapat diperluas dengan kasus error yang lebih banyak, termasuk nested procedure/function, array bertingkat, record di dalam array, dan kombinasi scope yang lebih kompleks.
